
Penelitian ini dilaksanakan dalam tujuh tahap utama sebagai berikut:

\begin{enumerate}
  \item \textbf{Studi Literatur.} Tahap awal ini akan fokus pada
  pengumpulan dan pemahaman mendalam teori dasar \emph{layered
  costmap}. Selain itu, akan dipelajari berbagai algoritma navigasi,
  seperti A*, yang relevan untuk \emph{path planning} robot, serta
  pemahaman mengenai teknik lokalisasi, seperti ICP. Tujuannya adalah
  untuk membangun dasar teoritis yang kuat sebagai landasan perancangan
  dan implementasi sistem navigasi.

  \item \textbf{Perancangan Sistem.} Pada tahap ini, struktur
  \emph{layered costmap} akan didesain secara rinci, termasuk
  penentuan jumlah lapisan dan jenis informasi yang akan dimuat pada
  setiap lapisan. Alur data dari berbagai sensor akan dirancang untuk
  memastikan informasi lingkungan dapat diintegrasikan dengan efisien
  ke dalam \emph{costmap}. Arsitektur sistem navigasi secara
  keseluruhan akan dibentuk, mempertimbangkan interaksi antara modul
  sensor LiDAR, lokalisasi, perencanaan jalur, dan kontrol gerak
  robot.

  \item \textbf{Perancangan Lingkungan Simulasi.} Lingkungan simulasi
  akan dirancang untuk merepresentasikan kondisi pasca gempa dalam
  skala laboratorium. Arena simulasi diadaptasi dari arena Kontes
  Robot SAR Indonesia (KRSRI) 2024 yang mencakup beberapa elevasi,
  tanjakan, dan permukaan tidak rata. Lingkungan ini dimodelkan di
  dalam simulator MuJoCo untuk keperluan pengujian sistem navigasi.

  \item \textbf{Implementasi Navigasi.} Modul \emph{layered costmap}
  yang telah dirancang akan diimplementasikan ke dalam sistem
  perangkat lunak robot, memastikan integrasi yang mulus dengan data
  sensor. Algoritma lokalisasi akan dikodekan dan disesuaikan untuk
  bekerja secara efektif dengan \emph{layered costmap}, memungkinkan
  robot mengetahui posisinya secara presisi di berbagai level.
  Terakhir, algoritma navigasi yang dipilih akan diintegrasikan untuk
  memungkinkan robot merencanakan dan mengeksekusi pergerakan antar
  titik tujuan, termasuk transisi antar elevasi.

  \item \textbf{Uji Coba dan Evaluasi.} Pengujian akan dilakukan untuk
  menguji kemampuan robot dalam bernavigasi antar titik tujuan yang
  terletak pada level berbeda di lingkungan simulasi. Pengujian ini
  akan mencakup skenario yang bervariasi, seperti navigasi melalui
  rintangan, menaiki dan menuruni elevasi, dan mencapai tujuan di
  region yang berbeda. Performa robot akan dievaluasi berdasarkan
  metrik seperti akurasi pencapaian tujuan, waktu tempuh, dan
  kelancaran pergerakan.

  \item \textbf{Analisis Hasil.} Data yang terkumpul dari tahap uji
  coba akan dianalisis secara mendalam untuk mengidentifikasi kekuatan
  dan kelemahan dari pendekatan navigasi yang menggunakan \emph{layered
  costmap}. Analisis juga akan mencakup identifikasi faktor-faktor
  yang mempengaruhi performa robot dan potensi perbaikan di masa
  mendatang.

  \item \textbf{Penyusunan Laporan.} Seluruh hasil penelitian, mulai
  dari studi literatur, perancangan sistem, implementasi, hingga
  analisis hasil uji coba, akan didokumentasikan secara lengkap dalam
  bentuk laporan. Laporan ini akan disusun sesuai dengan standar
  penulisan ilmiah sebagai tugas akhir.
\end{enumerate}

\newcommand{\w}{}
\newcommand{\G}{\cellcolor{gray}}
\begin{table}[H]
  \captionof{table}{Tabel timeline}
  \label{tbl:timeline}
  \begin{tabular}{|p{3.5cm}|c|c|c|c|c|c|c|c|c|c|c|c|c|c|c|c|}

    \hline
    \multirow{2}{*}{Kegiatan} & \multicolumn{16}{|c|}{Minggu}                                                                       \\
    \cline{2-17}              &
    1                         & 2                             & 3  & 4  & 5  & 6  & 7  & 8  & 9  & 10 & 11 & 12 & 13 & 14 & 15 & 16 \\
    \hline

    Studi Literatur          &
    \G                        & \G                            & \w & \w & \w & \w & \w & \w & \w & \w & \w & \w & \w & \w & \w & \w \\
    \hline

    Perancangan Sistem           &
    \w                        & \w                            & \G & \G & \G & \G & \w & \w & \w & \w & \w & \w & \w & \w & \w & \w \\
    \hline

    Perancangan Lingkungan Simulasi              &
    \w                        & \w                            & \w & \w & \w & \G & \G & \G & \w & \w & \w & \w & \w & \w & \w & \w \\
    \hline

    Implementasi Navigasi       &
    \w                        & \w                            & \w & \w & \w & \w & \w & \G & \G & \G & \G & \w & \w & \w & \w & \w \\
    \hline
    Uji Coba dan Evaluasi       &
    \w                        & \w                            & \w & \w & \w & \w & \w & \w & \w & \w & \G & \G & \G & \w & \w & \w \\
    \hline
    Analisis Hasil       &
    \w                        & \w                            & \w & \w & \w & \w & \w & \w & \w & \w & \w & \w & \w & \G & \w & \w \\
    \hline

     Penyusunan Laporan      &
    \w                        & \w                            & \w & \w & \w & \w & \w & \w & \w & \w & \w & \w & \w & \w & \G & \G \\
    \hline
  \end{tabular}
\end{table}
